\section{Primera vez como CEO: valor emocional, Ego y mentalidad de jugador de ajedrez}

Esta sección pasa del producto a los mecanismos del fundador. Ming Chaoping (明超平) es CEO por primera vez: disfruta del financiamiento y de la atención, pero también siente que camina sobre hielo delgado. La ventana para las aplicaciones de IA es muy corta y la forma de los productos cambia con rapidez; el equipo necesita estabilidad emocional, y el fundador necesita combatir su Ego de manera constante. No se trata de un chisme sobre la personalidad de alguien, sino de cómo una organización emprendedora mantiene la calidad de su juicio en condiciones de incertidumbre extrema.

\begin{figure}[H]
\centering
\includegraphics[width=0.94\textwidth]{ego-countermeasure.png}
\caption{Combatir el Ego: quien es CEO por primera vez debe contenerse a sí mismo con una perspectiva de tercero, valor emocional y mentalidad de jugador de ajedrez. Diagrama conceptual propio, elaborado a partir de la conversación de 02:14:50--02:36:57.}
\end{figure}

\begin{knowledgebox}{Lectura de la figura: el mecanismo del CEO también es un mecanismo de producto}
La perspectiva de tercero proviene del debate; el valor emocional, de la gestión de la organización; caminar sobre hielo delgado, de la conciencia del riesgo; combatir el Ego, de la autocalibración; y la mentalidad de jugador de ajedrez, del análisis de las partidas y del juicio sobre el orden de las jugadas. En conjunto, determinan si el equipo puede seguir tomando las decisiones correctas.
\end{knowledgebox}

\subsection{El usuario por encima del Ego}

Ming Chaoping menciona una frase que circula en el equipo: el usuario está por encima del Ego. Cada quien tiene sus propias opiniones y también interpreta los datos a su manera; pero la verdadera calibración proviene de hablar con los usuarios, estar inmerso en Discord y observar cómo crean, cómo comparten y cómo comentan. Para una empresa de aplicaciones de IA, la retroalimentación de los usuarios es especialmente importante, porque los escenarios nuevos no tienen métricas consolidadas y muchas veces la intuición solo puede construirse a partir del uso real.

\begin{warningbox}{El Ego se disfraza de juicio estratégico}
La confianza del fundador es útil, pero también puede llevarlo a confundir su gusto personal, el éxito en el financiamiento y los elogios externos con necesidades de los usuarios. Que el usuario esté por encima del Ego no significa no tener opiniones, sino que las opiniones deben calibrarse continuamente con los escenarios de los usuarios.
\end{warningbox}

\subsection{El jugador y su contrincante}

La metáfora del «ajedrez» al final de la entrevista es fundamental. Ming Chaoping dice que lo importante no es reproducir el registro final de la partida, sino entender por qué cada jugada se hizo donde se hizo. Con los productos ocurre lo mismo: ver hoy los productos terminados de WeChat, Douyin, Instagram y TikTok no significa que copiar su forma baste para tener éxito. Lo decisivo son el orden, las prioridades y el momento: qué se hace ahora, qué se hace dentro de tres años y qué no se hace nunca.

\begin{importantbox}{Las prioridades del producto son el orden de las jugadas}
Si la jugada 50 es un error, quizá no haya jugada 100. Con el emprendimiento de aplicaciones de IA sucede igual: empezar por el contenedor de creación, la distribución y la red de Agents, o empezar por los servicios empresariales, las suscripciones y las capacidades del modelo, cambia todo el camino posterior.
\end{importantbox}

\subsection{Resumen de la sección}

La dificultad de ser CEO por primera vez no se reduce al financiamiento y la contratación: consiste en calibrarse a sí mismo de manera continua. La segunda mitad de EP101 muestra que emprender con aplicaciones de Agents exige que el juicio de producto, el estado emocional de la organización, la calibración con los usuarios y el combate contra el Ego operen al mismo tiempo.
